En la ingeniería moderna, la teoría de conjuntos y la lógica proposicional no se resuelven solo con lápiz y papel. Los lenguajes de programación, como Python, integran estos conceptos en su núcleo. En esta sección aprenderemos a usar Python para operar conjuntos y evaluar proposiciones lógicas.

\subsection{Lógica Proposicional en Python}

Python utiliza el tipo de dato \texttt{bool} para manejar valores de verdad: \texttt{True} (Verdadero) y \texttt{False} (Falso). Los conectores lógicos que hemos estudiado tienen su equivalente directo en palabras reservadas del lenguaje.

\begin{table}[ht]
\caption{Equivalencia Lógica - Python}\label{tabla:python_logic}
\centering
\begin{tabular}{|l|c|c|l|}
\hline
\textbf{Operación} & \textbf{Símbolo} & \textbf{En Python} & \textbf{Ejemplo} \\ \hline
Negación & $\lnot p$ & \texttt{not p} & \texttt{not True} $\to$ \texttt{False} \\ \hline
Conjunción & $p \land q$ & \texttt{p and q} & \texttt{True and False} $\to$ \texttt{False} \\ \hline
Disyunción & $p \lor q$ & \texttt{p or q} & \texttt{True or False} $\to$ \texttt{True} \\ \hline
Disyunción Excl. & $p \veebar q$ & \texttt{p \^{} q} & \texttt{True \^{} True} $\to$ \texttt{False} \\ \hline
Implicación & $p \Rightarrow q$ & \texttt{not p or q} & (No tiene operador nativo) \\ \hline
\end{tabular}
\end{table}

\noindent \textbf{Ejemplo 1: Evaluación de una Tautología}
Queremos comprobar si la expresión $(p \land q) \Rightarrow p$ es siempre verdadera.

\begin{lstlisting}[caption=Comprobación de Tautología]
# Definimos una funcion para la implicacion
def implicacion(p, q):
    return not p or q

print("p \t q \t Resultado")
print("-" * 25)

# Iteramos sobre todos los valores posibles (Tabla de Verdad)
for p in [True, False]:
    for q in [True, False]:
        antecedente = p and q
        resultado = implicacion(antecedente, p)
        print(f"{p} \t {q} \t {resultado}")
\end{lstlisting}

\noindent \textit{Salida del programa:}
\begin{lstlisting}[language=bash, numbers=none, frame=none, backgroundcolor=\color{white}]
p 	     q 	     Resultado
-------------------------
True 	 True 	 True
True 	 False 	 True
False 	 True 	 True
False 	 False 	 True
\end{lstlisting}
Como vemos, el resultado es siempre \texttt{True}, confirmando que es una tautología.

\subsection{Teoría de Conjuntos en Python}

Python tiene una estructura de datos nativa llamada \texttt{set} (conjunto), que cumple todas las propiedades matemáticas: no tiene elementos repetidos y no tiene orden.

\begin{table}[ht]
\caption{Operaciones de Conjuntos en Python}\label{tabla:python_sets}
\centering
\begin{tabular}{|l|c|l|}
\hline
\textbf{Operación} & \textbf{Símbolo} & \textbf{Sintaxis en Python} \\ \hline
Definición & $\{1, 2\}$ & \texttt{A = \{1, 2\}} \\ \hline
Unión & $A \cup B$ & \texttt{A | B}  o  \texttt{A.union(B)} \\ \hline
Intersección & $A \cap B$ & \texttt{A \& B}  o  \texttt{A.intersection(B)} \\ \hline
Diferencia & $A - B$ & \texttt{A - B}  o  \texttt{A.difference(B)} \\ \hline
Diferencia Sim. & $A \Delta B$ & \texttt{A \^{} B}  o  \texttt{A.symmetric\_difference(B)} \\ \hline
Subconjunto & $A \subseteq B$ & \texttt{A <= B}  o  \texttt{A.issubset(B)} \\ \hline
\end{tabular}
\end{table}

\noindent \textbf{Ejemplo 2: Operaciones básicas}
Dados $A = \{1, 2, 3, 4\}$ y $B = \{3, 4, 5, 6\}$, calcular su intersección y diferencia simétrica.

\begin{lstlisting}[caption=Operaciones de Conjuntos]
A = {1, 2, 3, 4}
B = {3, 4, 5, 6}

interseccion = A & B
dif_simetrica = A ^ B

print(f"A interseccion B = {interseccion}")
print(f"A Delta B = {dif_simetrica}")
\end{lstlisting}

\noindent \textit{Salida:}
\begin{lstlisting}[language=bash, numbers=none, frame=none, backgroundcolor=\color{white}]
A interseccion B = {3, 4}
A Delta B = {1, 2, 5, 6}
\end{lstlisting}

\subsection{Ejercicios Resueltos}

\textbf{Ejercicio 1: Comprobación de Leyes de De Morgan}
Demuestre usando Python que $\lnot (p \lor q) \iff (\lnot p \land \lnot q)$.

\textit{Solución:}
\begin{lstlisting}[caption=Validación Ley de De Morgan]
valores = [True, False]
es_equivalente = True

for p in valores:
    for q in valores:
        lado_izq = not (p or q)
        lado_der = (not p) and (not q)
        
        if lado_izq != lado_der:
            es_equivalente = False
            break

if es_equivalente:
    print("La ley de De Morgan se cumple para todos los casos.")
else:
    print("La ley no se cumple.")
\end{lstlisting}

\textbf{Ejercicio 2: Filtrado por Comprensión}
Dado el conjunto universal $U = \{0, 1, 2, ..., 20\}$, cree un conjunto $A$ con los números pares y un conjunto $B$ con los múltiplos de 3. Encuentre $A \cap B$.

\textit{Solución:}
\begin{lstlisting}[caption=Conjuntos por Comprensión]
# Conjunto Universal (range crea una secuencia de 0 a 20)
U = set(range(21)) 

# Definicion por comprension (Set Comprehension)
# Sintaxis: {x for x in U if condicion}
A = {x for x in U if x % 2 == 0}
B = {x for x in U if x % 3 == 0}

resultado = A & B
print(f"Numeros pares y multiplos de 3: {resultado}")
\end{lstlisting}

\subsection{Ejercicios Propuestos con Python}

Para resolver en el computador:

\begin{enumerate}
    \item \textbf{Generador de Tablas de Verdad:} Escriba un programa en Python que pida al usuario el valor de tres variables booleanas ($p, q, r$) e imprima el valor de verdad de la expresión $(p \lor q) \land (\lnot r)$.
    
    \item \textbf{Validador de Subconjuntos:} Defina dos listas de números ingresadas por teclado, conviértalas a conjuntos y determine si la primera es subconjunto de la segunda utilizando el operador \texttt{<=}.
    
    \item \textbf{Diferencia de Listas:} Suponga que tiene una lista de estudiantes inscritos en el curso de ``Lógica'' y otra lista de estudiantes en ``Matemáticas''. Use Python para encontrar qué estudiantes están viendo ``Lógica'' pero \textbf{no} ``Matemáticas'' (Diferencia de conjuntos $L - M$).
    
    \item \textbf{El Problema de los 24 Libros:} Retome el Ejercicio 9 de la sección de conjuntos. Escriba un script que defina los cardinales de cada intersección y use la fórmula del Principio de Inclusión-Exclusión para calcular la unión de los tres conjuntos ($|A \cup B \cup C|$). Luego reste ese valor al total (24) para hallar cuántos libros no tratan ninguno de los temas.
\end{enumerate}
